\section{可微映射}\label{sec:differentiable-maps}

在本书中, \textbf{可微}和\textbf{光滑}都指\textbf{无穷次可微} ($C^\infty$).

\begin{RiemannDefinition}\label{def:differentiable-map}
设 $M$ 和 $N$ 分别为 $m$ 维和 $n$ 维可微流形. 若存在 $M$ 在 $p\in M$ 处的参数化 $(U,\varphi)$ (即 $p\in\varphi(U)$) 及 $N$ 在 $f(p)$ 处的参数化 $(V,\psi)$, 满足 $f(\varphi(U))\subset\psi(V)$, 且映射
\[
  \widehat f:=\psi^{-1}\circ f\circ\varphi:U\subset\mathbb{R}^m\to\mathbb{R}^n
\]
光滑 (Fig.~\ref{fig:local-representation-of-manifold-map}), 则称映射 $f:M\to N$ 在 $p$ 处\textbf{可微} (或\textbf{光滑}, 或为 \textbf{$C^\infty$}).

若 $f$ 在 $M$ 的某个子集的每一点都可微, 就称 $f$ 在该子集上\textbf{可微}.
\end{RiemannDefinition}

\RGfigure[.56\textwidth]{local-representation-of-manifold-map}{流形之间映射的局部表示}

因为坐标变换是光滑的, 这个定义与在 $f(p)$ 和 $p$ 处选取的参数化无关. 映射 $\widehat f:=\psi^{-1}\circ f\circ\varphi:U\subset\mathbb{R}^m\to\mathbb{R}^n$ 称为 $f$ 的\textbf{局部表示}, 即 $f$ 在由 $\varphi$ 和 $\psi$ 确定的局部坐标下的表达式. 一切光滑函数 $f:M\to N$ 的集合记为 $C^\infty(M,N)$, 并把 $C^\infty(M,\mathbb{R})$ 简记为 $C^\infty(M)$.

两个流形之间的可微映射 $f:M\to N$ 是连续的 [参见 Exercise~\ref{ex:properties-of-differentiable-maps}(\ref{ex:smooth-map-is-continuous})]. 此外, 若 $f$ 是双射且其逆映射 $f^{-1}:N\to M$ 也可微, 则称 $f$ 为\textbf{微分同胚}. 若可微流形 $M$ 和 $N$ \textbf{微分同胚}, 即存在微分同胚 $f:M\to N$, 就把它们视为相同. 如果存在 $p\in M$ 的邻域 $V$ 和 $f(p)$ 的邻域 $W$, 使得 $f|_V:V\to W$ 是微分同胚, 则称 $f$ 在 $p$ 处是\textbf{局部微分同胚}.

很长一段时间里, 人们认为每个拓扑流形在相差微分同胚的意义下只有一种可微结构 (Exercise~\ref{ex:exercise-differentiable-structures}(\ref{ex:prove-inequivalence-of-line-atlases}) 和 Exercise~\ref{ex:properties-of-differentiable-maps}(\ref{ex:properties-of-differentiable-maps-part-six}) 中的两个不同可微结构是微分同胚的, 参见 Exercise~\ref{ex:properties-of-differentiable-maps}(\ref{ex:properties-of-differentiable-maps-part-six})). 但在 1956 年, Milnor \RGcite{Mil56} 给出了与 $S^7$ 同胚但不微分同胚的流形的例子. 后来, Milnor 和 Kervaire [\RGref{bib:Mil59}, \RGref{bib:KM63}] 证明了还有一些维数大于 7 的球面也容许多种可微结构. 例如, $S^{19}$ 有 73 种不同的光滑结构, $S^{31}$ 有 16,931,177 种. 时间更近一些, Freedman \RGcite{Fre82} 和 Gompf \RGcite{Gom83} 分别在 1982 年和 1983 年构造了 $\mathbb{R}^4$ 上非标准可微结构的例子.

\begin{RiemannExercise}\label{ex:properties-of-differentiable-maps}
\begin{enumerate}[label=(\arabic*),ref=\arabic*,leftmargin=*]
\item\label{ex:independence-from-parametrization} 证明 Definition~\ref{def:differentiable-map} 不依赖于参数化的选取.
\item\label{ex:smooth-map-is-continuous} 证明两个光滑流形之间的可微映射 $f:M\to N$ 是连续的.
\item\label{ex:composite-of-smooth-maps} 证明若 $f:M_1\to M_2$ 和 $g:M_2\to M_3$ 是光滑流形 $M_1,M_2,M_3$ 之间的可微映射, 则 $g\circ f:M_1\to M_3$ 也可微.
\item\label{ex:smooth-antipodal-map} 证明由 $f(x)=-x$ 给出的\textbf{对径映射} $f:S^n\to S^n$ 是可微的.
\item\label{ex:mobius-transformations-of-the-sphere} 利用从北极出发的球极投影 $\pi_N:S^2\setminus\{N\}\to\mathbb{R}^2$, 并把 $\mathbb{R}^2$ 与复平面 $\mathbb{C}$ 等同, 可以把 $S^2$ 与 $\mathbb{C}\cup\{\infty\}$ 等同, 其中 $\infty$ 称为\textbf{无穷远点}. \textbf{Möbius 变换}是形如
\[
  f(z)=\frac{az+b}{cz+d},
\]
的映射 $f:\mathbb{C}\cup\{\infty\}\to\mathbb{C}\cup\{\infty\}$, 其中 $a,b,c,d\in\mathbb{C}$ 满足 $ad-bc\ne 0$, 且 $\infty$ 满足
\[
  \frac{\alpha}{\infty}=0,\qquad \frac{\alpha}{0}=\infty
\]
对任意 $\alpha\in\mathbb{C}\setminus\{0\}$ 成立. 证明任意 Möbius 变换 $f$ 视为映射 $f:S^2\to S^2$ 时都是微分同胚. {\small (\textbf{Hint:} 先证明任意 Möbius 变换都是形如 $g(z)=\frac{1}{z}$ 和 $h(z)=az+b$ 的变换的复合.)}
\item\label{ex:properties-of-differentiable-maps-part-six} 再考虑 Example~\ref{ex:examples-of-differentiable-structures}(\ref{ex:inequivalent-atlases-on-the-line}) 和 Exercise~\ref{ex:exercise-differentiable-structures}(\ref{ex:prove-inequivalence-of-line-atlases}) 中 $\mathbb{R}$ 上的两个图册 $\{(\mathbb{R},\varphi_1)\}$ 和 $\{(\mathbb{R},\varphi_2)\}$, 其中 $\varphi_1(x)=x$, $\varphi_2(x)=x^3$. 证明:
\begin{enumerate}[label=(\alph*),leftmargin=*]
\item 恒等映射 $i:(\mathbb{R},\varphi_1)\to(\mathbb{R},\varphi_2)$ 不是微分同胚;
\item 由 $f(x)=x^3$ 给出的映射 $f:(\mathbb{R},\varphi_1)\to(\mathbb{R},\varphi_2)$ 是微分同胚 (这说明这两个图册虽然定义了不同的可微结构, 但这些结构彼此微分同胚).
\end{enumerate}
\end{enumerate}
\end{RiemannExercise}
